\section*{Chapitre \ref{ch:g11:titration} --- Le titrage}

\begin{solution}{exo:g11:titration:1}
La \omterm{def:g11:titration:titration}{solution titrante} est la solution de concentration connue avec
précision, dans la burette~; la \omterm{def:g11:titration:titration}{solution titrée} contient l'espèce à
doser, dans l'erlenmeyer. À l'\omterm{def:g11:titration:equivalence}{équivalence}, le \omterm{def:g7:chemical-reactions:reactant}{réactif} titrant ajouté et
l'espèce titrée sont dans les proportions de l'équation~: tous deux
sont entièrement consommés.
\end{solution}

\begin{solution}{exo:g11:titration:2}
$n = 0.0200 \times 0.0125 = \qty{2.50e-4}{mol}$.
\end{solution}

\begin{solution}{exo:g11:titration:3}
$V_E = \qty{14.3}{mL}$~; la droite du fer(II) part de
\qty{14.3e-4}{mol}, soit \qty{1.43e-3}{mol}.
\end{solution}

\begin{solution}{exo:g11:titration:4}
Le \omterm{def:g7:chemical-reactions:reactant}{réactif} titrant est lui-même coloré (violet) et son produit presque
incolore~: la couleur apparaît dans l'erlenmeyer dès que le permanganate
est en excès.
\end{solution}

\begin{solution}{exo:g11:titration:5}
5~: l'équation consomme 5 \omterm{def:g9:ions:ion}{ions} fer(II) par \omterm{def:g9:ions:ion}{ion} permanganate, et à
l'\omterm{def:g11:titration:equivalence}{équivalence} les \omterm{def:g7:chemical-reactions:reactant}{réactifs} ont été réunis dans ces proportions.
\end{solution}

\begin{solution}{exo:g11:titration:6}
$n(\ce{MnO4^-}) = 0.0200 \times 0.00840 = \qty{1.68e-4}{mol}$~;
$n(\ce{Fe^{2+}}) = 5 \times \num{1.68e-4} = \qty{8.40e-4}{mol}$~;
$c = \num{8.40e-4} / 0.0100 = \qty{0.0840}{mol/L}$~;
$C_m = 0.0840 \times 55.8 = \qty{4.69}{g/L}$.
\end{solution}

\begin{solution}{exo:g11:titration:7}
Rapport 1/1~: $n = 0.0100 \times 0.0114 = \qty{1.14e-4}{mol}$ dans
\qty{10.0}{mL}, donc \qty{1.14e-3}{mol} dans les \qty{100.0}{mL}, soit
$\num{1.14e-3} \times 176.0 = \qty{0.201}{g}$~: environ \qty{200}{mg}.
\end{solution}

\begin{solution}{exo:g11:titration:8}
Chaque goutte ne réagirait pas aussitôt~: la couleur de la \omterm{def:g11:titration:titration}{solution
titrante} persisterait avant l'\omterm{def:g11:titration:equivalence}{équivalence}, et on croirait voir la fin
trop tôt (ou il faudrait attendre plusieurs minutes après chaque
goutte).
\end{solution}

\begin{solution}{exo:g11:titration:9}
Trop grand~: elle a dépassé l'\omterm{def:g11:titration:equivalence}{équivalence} en ajoutant du permanganate en
excès. Près de l'\omterm{def:g11:titration:equivalence}{équivalence}, elle aurait dû ajouter la \omterm{def:g11:titration:titration}{solution
titrante} goutte à goutte, en agitant, et s'arrêter au premier rose pâle
persistant.
\end{solution}

\begin{solution}{exo:g11:titration:10}
$n(\ce{Fe^{2+}}) = 5 \times 0.0200 \times 0.0130 = \qty{1.30e-3}{mol}$
dans \qty{10.0}{mL}~: $c = \qty{0.130}{mol/L}$ pour la solution diluée,
et $10 \times 0.130 = \qty{1.30}{mol/L}$ pour la solution concentrée.
\end{solution}

\begin{solution}{exo:g11:titration:11}
$0.0630 \times 55.8 = \qty{3.52}{g}$.
\end{solution}

\begin{solution}{exo:g11:titration:12}
$0.1 / 7.2 \approx \qty{1.4}{\percent}$, contre \qty{0.7}{\percent} pour
\qty{14.3}{mL}. Titrer un plus grand volume (ou un échantillon plus
concentré) de la solution, ou utiliser une \omterm{def:g11:titration:titration}{solution titrante} plus
diluée.
\end{solution}

\begin{solution}{exo:g11:titration:13}
$n(\ce{MnO4^-}) = 0.0200 \times 0.0176 = \qty{3.52e-4}{mol}$~;
$n(\ce{H2O2}) = \frac{5}{2} \times \num{3.52e-4} = \qty{8.80e-4}{mol}$
dans \qty{10.0}{mL}~: \qty{0.0880}{mol/L} dans la solution diluée,
\qty{0.880}{mol/L} dans le désinfectant, soit
$0.880 \times 34.0 = \qty{29.9}{g/L}$ (une solution à 3~\% environ).
\end{solution}

\begin{solution}{exo:g11:titration:14}
$n(\ce{MnO4^-}) = \num{1.4e-3} / 5 = \qty{2.8e-4}{mol}$, à délivrer en
environ \qty{0.015}{L}~: $c \approx \qty{0.019}{mol/L}$, d'où une
solution à \qty{0.0200}{mol/L}. Dix fois plus concentrée, elle donnerait
$V_E$ voisin de \qty{1.4}{mL}~: l'erreur de la burette atteindrait
environ \qty{7}{\percent}.
\end{solution}

\begin{solution}{exo:g11:titration:15}
$8 \times 0.0200 \times 0.0143 = \qty{2.29e-3}{mol}$ d'\omterm{def:g9:ions:ion}{ions} hydrogène.
\qty{10}{mL} à \qty{1.0}{mol/L} en contiennent \qty{1.0e-2}{mol}~:
c'est assez, plus de quatre fois. Sans acide, la réaction du chapitre ne
pourrait pas se produire telle qu'elle est écrite (les \omterm{def:g9:ions:ion}{ions} hydrogène
sont un \omterm{def:g7:chemical-reactions:reactant}{réactif}), et le permanganate réagirait autrement.
\end{solution}

\begin{solution}{pb:g11:titration:1}
\textbf{1.} \ce{MnO4^-}/\ce{Mn^{2+}} et \ce{Fe^{3+}}/\ce{Fe^{2+}}.
L'\omterm{def:g11:redox:oxidant}{oxydant} est l'\omterm{def:g9:ions:ion}{ion} permanganate, le \omterm{def:g11:redox:oxidant}{réducteur} l'\omterm{def:g9:ions:ion}{ion} fer(II).

\textbf{2.} \ce{MnO4^- + 8H+ + 5e- -> Mn^{2+} + 4H2O}~;
\ce{Fe^{2+} -> Fe^{3+} + e-}.

\textbf{3.} On multiplie la seconde par 5~:
\ce{MnO4^- + 8H+ + 5Fe^{2+} -> Mn^{2+} + 5Fe^{3+} + 4H2O}. Charge à
gauche~: $-1 + 8 + 10 = +17$~; à droite~: $+2 + 15 = +17$.

\textbf{4.} Les \omterm{def:g9:ions:ion}{ions} hydrogène sont un \omterm{def:g7:chemical-reactions:reactant}{réactif}~: la réaction a besoin
d'un milieu acide.

\textbf{5.} Elle est totale, rapide, et c'est la seule réaction du
permanganate dans l'erlenmeyer~; le permanganate, violet, est consommé
avant l'\omterm{def:g11:titration:equivalence}{équivalence} et colore l'erlenmeyer en rose après elle.

\textbf{6.} Dans la burette~; on la lit au bas du ménisque, l'œil à sa
hauteur, avant et après, à \qty{0.05}{mL} près.

\textbf{7.} Jaune pâle (\omterm{def:g9:ions:ion}{ions} fer)~; chaque goutte violette est aussitôt
consommée par les \omterm{def:g9:ions:ion}{ions} fer(II) encore présents.

\textbf{8.} $0.0200 \times 0.0143 = \qty{2.86e-4}{mol}$.

\textbf{9.} $n(\ce{Fe^{2+}}) = 5 \times \num{2.86e-4} =
\qty{1.43e-3}{mol}$.

\textbf{10.} $\num{1.43e-3} \times 55.8 = \qty{0.0798}{g} =
\qty{79.8}{mg}$.

\textbf{11.} $(80 - 79.8)/80 \approx \qty{0.3}{\percent}$.

\textbf{12.} $\qty{0.1}{mL}$ sur \qty{14.3}{mL}, c'est \qty{0.7}{\percent},
soit environ \qty{0.6}{mg}~: le résultat, $79.8 \pm 0.6$~mg, est en
accord avec l'étiquette.

\textbf{13.} $14.1$~: $5 \times 0.0200 \times 0.0141 \times 55.8 =
\qty{78.7}{mg}$~; $14.4$~: \qty{80.4}{mg}.

\textbf{14.} $(79.8 + 78.7 + 80.4)/3 = \qty{79.6}{mg}$.

\textbf{15.} Les comprimés eux-mêmes diffèrent légèrement par leur
teneur en fer, d'environ \qty{1}{\percent}.

\textbf{16.} Non~: $V_E$ vaudrait environ \qty{2.9}{mL}, et l'erreur de
la burette \qty{3.5}{\percent}, cinq fois plus.

\textbf{17.} La vitamine C serait elle aussi oxydée par le
permanganate~: la réaction ne serait plus unique, $V_E$ serait trop
grand, et le fer serait surestimé.

\textbf{18.} $\num{1.43e-3} \times 151.9 = \qty{0.217}{g}$, environ
\qty{217}{mg}.

\textbf{19.} $\num{1.43e-3} \times \num{6.02e23} = \num{8.6e20}$ \omterm{def:g9:ions:ion}{ions}.

\textbf{20.} Environ \qty{80}{mg} de fer par comprimé~: l'étiquette est
honnête.
\end{solution}
